% !TEX root = ../main.tex
\chapter{Lagrangians and action functionals (Lagrangefunktionen und Wirkungsfunktionale)}\label{ch:c13}
\begin{paracol}{2}

\begin{sources}
\begin{tabularx}{\linewidth}{@{}l l L@{}}
\toprule
\textbf{Video} & \textbf{Length} & \textbf{Content}\\
\midrule
\seg{13}{1}{Lecture 13: Lagrangians and action functionals} & 1:44:09 & Start of part C, the Lagrangian formulation: motivation (a complicated system with constraints); the tangent bundle and the bundle projection; the induced atlas and its smoothness; lifting of curves and the velocity phase space; the Lagrangian as a smooth function on the tangent bundle; its chart representation, the notation of the textbooks, the harmonic oscillator and its level lines; an outlook on the Hamiltonian formulation; the action functional; the Feynman path integral as an outlook; goal and strategy of the calculus of variations; modelling in physics.\\
\bottomrule
\end{tabularx}

\smallskip
\textbf{Transcript} \file{materials/transcripts/C13-V01.txt} (German YouTube auto-captions, mediocre in the spoken formulas; the formulas are reconstructed by recomputing and marked [?] where unsure).
\end{sources}

\begin{reading}[Where this lecture sits in the other notes]
\emph{Gravity and Light}: lecture 14, ``Matter'' (the Lagrangian part: Lagrangians and the action there). \emph{Geometric Anatomy}: none.
Continues \cref{ch:c12} (the two-body problem, the end of the Newtonian part B) and uses charts, tangent spaces and the differential of \cref{ch:c02,ch:c03,ch:c04}; the next lecture develops the calculus of variations.
\end{reading}

% ==================================================================
\section{Part C: the Lagrangian formulation of mechanics}\label{sec:lag-intro}

Good morning, please take your seats. Today a new part of the lecture starts, \textbf{part C: the Lagrangian formulation of mechanics (Lagrange-Formulierung der Mechanik)} \ts{13}{1}{0:00}. So that we do not lose sight of what we want: we want to understand how such a system moves when I give it a push, and above all we want to set up the equations of motion for such a system \emph{without} thinking very much.

Look at the system on the handout (the drawing of the lecture): it shows two things. First, it is a rather \emph{complex system}, far from the simple flag-swinger (``Fahnenschwinger''), which one could get under control; here are various cross-couplings between the masses, there are rods. Second, there are \textbf{constraints} (Zwangsbedingungen): this bead can move only along this bent piece of wire, it is forced to be on it; this mass is forced to be wherever the rigid pendulum rod allows it;
this other mass is even more complicated, since where it may be depends on where that thing is; and down here there is a rather free mass that swings back and forth but can also move in the other direction: ``quite a chaos'' \ts{13}{1}{0:21}--\ts{13}{1}{1:45}.
It sounds as if one wanted too much at once: to understand such a thing, complex and with constraints. But the \emph{laconic formulation of mechanics}, which really is only a mathematical reformulation of Newton's law, kills both problems at once, and it is child's play to set up the equation of motion. Instead of modelling \emph{forces} one models no vector field or interaction field,
but a \emph{single scalar function}, the so-called \textbf{Lagrangian (Lagrange-Funktion)} describing the system; this I will show you in the course of the next lectures. What matters: this Lagrangian is not a scalar function on space but a scalar function on a space of \emph{twice the dimension}, the so-called \textbf{velocity phase space} (Geschwindigkeitsphasenraum), one could also say the space of initial conditions \ts{13}{1}{1:45}--\ts{13}{1}{2:48}.
We will develop this mathematically and provide the necessary technology.

\begin{coursenote}[Books recommended (as said in class)]\label{cn:lag-books}
Three books are recommended: one by Scheck (``he does such things a bit more geometrically, as we do; we do it a bit differently again''), Flie{\ss}bach's \emph{Mechanik} as a standard work, and a third, (name garbled in the captions, ``the book of Bas''), with many nice examples [?] \ts{13}{1}{2:48}--\ts{13}{1}{3:29}.
\end{coursenote}

Today we begin with the development of the formal background, and the first concept to define is the \emph{tangent bundle} \ts{13}{1}{3:08}.

% ==================================================================
\section{The tangent bundle}\label{sec:lag-tb}

I first carry out some mathematical developments that are applicable quite arbitrarily; so I will always deal with a smooth manifold $M$ of dimension $d$, but you may keep in mind that in our application it will later be the \emph{absolute space} (to which a metric is added); you should not think of spacetime, which we also called $M$; here it is a general mathematical development; in the simplest case the space $\R^3$ \ts{13}{1}{3:53}--\ts{13}{1}{4:35}.
So we have a set $M$, and if I look first only at this set I can define a new set. Recall that $T_pM$ was the tangent space to the manifold at one point $p$ (\cref{def:ts-space}); the new set is called $TM$ and is built from \emph{all} the tangent spaces to a smooth manifold, by taking every single tangent space and gluing them together in one big disjoint union.

\begin{definition}[New concepts: the tangent bundle (Tangentialb\"undel) as a set, and the bundle projection]\label{def:lag-tb}
For a smooth manifold $(M,\mathcal O,\mathcal A)$ of dimension $d$ define the set
\begin{equation}\label{eq:lag-tb}
  TM\;:=\;\bigsqcup_{p\in M}T_pM
\end{equation}
(the disjoint union of all tangent spaces), and the \textbf{bundle projection} $\pi\colon TM\to M$ by
\begin{equation}\label{eq:lag-pi}
  \pi(X)\;:=\;\text{the unique point }p\in M\text{ with }X\in T_pM .
\end{equation}
The pair $(TM,\pi)$, the set $TM$ together with the map $\pi$, is called the \textbf{tangent bundle of $M$} \ts{13}{1}{3:53}--\ts{13}{1}{11:53}.
\end{definition}
(Why ``bundle'': the tangent spaces are bundled up, as one bundles a fascicle, see below.)

Explanation of $\pi$, as said in the lecture: take an element $X\in TM$. $TM$ is a union of tangent spaces, so $X$ lies in \emph{some} $T_pM$; since the union is disjoint, there is exactly one such $p$, and we give it the name $\pi(X)$, so that it need not always be written in a complicated way. So $\pi$ takes an element of the bundle and throws it back to a point of the original manifold: ``where is this vector attached''. As a set, $TM$ is nothing more than that; more structure will be added \ts{13}{1}{5:36}--\ts{13}{1}{6:39}.

\begin{example}[The circle: its tangent bundle is a cylinder]\label{ex:lag-cyl}
Take $M=S^1$. At each point $p$ there is the tangent space (``the tangent, it goes on as far as you like''); at another point another one; in the drawing they seem to cross each other, but of course the picture lies a little: the spaces have nothing to do with each other as sets, it is only the picture which makes them intersect. The set $TS^1$ is a ``pulverised set'' that we can grind in a mortar, put into a small box and reassemble: it is a \emph{cylinder}, as a set.
It does not stop at the top, because the tangent spaces are isomorphic to $\R$ (infinite lines): all points of the cylinder are elements of $TS^1$. If a vector lies at the point $p$, then as an element of the bundle it is a point on the cylinder above $p$, namely at the height given by its length: $X\in T_{\pi(X)}S^1$ \ts{13}{1}{6:39}--\ts{13}{1}{10:51}.
\end{example}
The statement ``there is a base point'' is thus read as: for each element $X$ of the big set I find a base point $p=\pi(X)$ such that $X$ lies in the tangent space over this point; this is where the vector is \emph{attached} \ts{13}{1}{10:09}--\ts{13}{1}{10:51}. $\pi$ is a ``wonderful object'' which brings an element of the bundle back to the manifold; $TM$ together with $\pi$ is the tangent bundle: ``simply all tangent spaces taken together'' \ts{13}{1}{10:51}--\ts{13}{1}{11:53}.

\subsection{The tangent bundle is a smooth manifold: the induced atlas}\label{ssec:lag-atlas}

So far, apart from the projection, only a new, big set. For our further investigations it would be great if this set, which we have constructed from a manifold, were itself a \emph{smooth manifold}. How do we do it? If $M$ is a smooth manifold, then we have a smooth atlas $\mathcal A$; we try to \emph{construct from the atlas of the base manifold an atlas of the whole thing}, and then examine (1) whether the constructed atlas is an atlas, and (2) whether it is smooth. If this succeeds, we have a smooth manifold (it will be the velocity phase space) \ts{13}{1}{11:53}--\ts{13}{1}{13:17}. The answer is yes, and we proceed as follows.

Consider a chart $(U,x)$ from the atlas of the manifold with which we begin (the lecturer draws a one-dimensional manifold $M$ with the tangent spaces at its points, and a chart domain $U$ and the chart map $x\colon U\to x(U)\subseteq\R^d$) \ts{13}{1}{13:41}--\ts{13}{1}{15:48}.
We want to make a chart of $TM$, of the whole object, and for the chart domain we take the tangent spaces over this open set only:
\begin{equation}\label{eq:lag-TU}
  TU\;:=\;\bigsqcup_{p\in U}T_pM\;=\;\pi^{-1}(U),
\end{equation}
``the collection of only these blue tangent spaces'': outside $U$ there is still a tangent space to the manifold, but its points do not belong to the chart domain $TU$. This is simply the bundle over the chart domain \ts{13}{1}{15:48}--\ts{13}{1}{16:50}. On it we construct a chart map $\xi_x$ (``the Greek letter $\xi$, the chart map which we shall construct from the old chart $x$''), from the old chart $x$ on the \emph{base manifold} (this is the name for $M$, as the bundle sits over it), which lands in $\R^{2d}$: ``so the new manifold has twice the dimension'' \ts{13}{1}{16:50}--\ts{13}{1}{18:14}.

\begin{definition}[New concepts: the induced chart $\xi_x$ on $TU$]\label{def:lag-xi}
For an element $X\in TU$ define
\begin{equation}\label{eq:lag-xi}
  \xi_x(X)\;:=\;\Bigl(x^1\bigl(\pi(X)\bigr),\dots,x^d\bigl(\pi(X)\bigr),\ (\mathrm dx^1)_{\pi(X)}(X),\dots,(\mathrm dx^d)_{\pi(X)}(X)\Bigr)\ \in\R^{2d},
\end{equation}
that is, the first $d$ components are the chart coordinates of the base point $\pi(X)$, and the last $d$ components are obtained by letting the differentials $(\mathrm dx^i)_{\pi(X)}$ (\cref{def:ts-differential}) at the base point act on $X$ \ts{13}{1}{18:14}--\ts{13}{1}{21:48}.
\end{definition}
First a check that this is sensible: $X\in TU$, so $\pi(X)\in U$, and the component functions $x^i$ can act on $\pi(X)$ and give real numbers; since $X$ is an element of $T_{\pi(X)}M$ the differentials, which are covectors at $\pi(X)$ (linear maps $T_{\pi(X)}M\to\R$), can act on $X$ and again give a real number. So we have constructed an element of $\R^{2d}$ \ts{13}{1}{19:39}--\ts{13}{1}{21:48}.

The first $d$ numbers are understood: take a point of the tangent bundle, project it with $\pi$ to the base manifold, and use for this point the normal chart: the coordinates of the base point. What are the last $d$ numbers? Since $X\in T_{\pi(X)}M$ we can expand it in the basis induced by the chart $x$ (\cref{def:ts-basis}):
\begin{equation}\label{eq:lag-comp}
  X\;=\;X^m\,\frac{\partial}{\partial x^m}\bigg|_{\pi(X)},\qquad (\mathrm dx^i)_{\pi(X)}(X)\;=\;X^m\,\delta^i_m\;=\;X^i ,
\end{equation}
by the duality of the dual basis ($(\mathrm dx^i)(\partial/\partial x^m)=\delta^i_m$): ``we have computed this a hundred times''; so the last $d$ numbers are just the \emph{components of the vector $X$} with respect to the basis induced by the same chart \ts{13}{1}{22:30}--\ts{13}{1}{25:39}. In the one-dimensional picture of the circle mapped to a chart: the circle is mapped to part of the real axis, and the tangent spaces are ``set perpendicular'': the horizontal axis gives the coordinate $x^1$ of the base point and the vertical the component of the vector; ``this is the chart image'' \ts{13}{1}{25:39}--\ts{13}{1}{26:21}.

\begin{definition}[New concepts: the induced atlas $\mathcal A_{TM}$]\label{def:lag-atlas}
The \textbf{induced atlas} $\mathcal A_{TM}$ of $TM$ consists of all the charts $(TU,\xi_x)$ constructed in this way from all charts $(U,x)$ of $\mathcal A$ \ts{13}{1}{26:42}--\ts{13}{1}{27:24}.
\end{definition}
It is easy to see that it covers everything (the $U$ cover $M$, hence the $TU$ cover $TM$), so it is an atlas. But \emph{is it smooth?} How does one examine whether an atlas is smooth? By looking at the chart transition maps (\cref{def:dm-transition}) and checking if they are differentiable \ts{13}{1}{27:24}--\ts{13}{1}{28:32}.

Take two charts $(TU,\xi_x)$ and $(TV,\xi_y)$ of the induced atlas, with $U\cap V\neq\emptyset$, constructed from charts $(U,x)$ and $(V,y)$ of the base atlas. The transition map goes from $\xi_x(TU\cap TV)$ to $\xi_y(TU\cap TV)$ (the same as $\xi(T(U\cap V))$) and, in components, an element of $\xi_x(TU\cap TV)\subseteq\R^{2d}$ is a $2d$-tuple $(\alpha^1,\dots,\alpha^d,\beta^1,\dots,\beta^d)$ \ts{13}{1}{28:32}--\ts{13}{1}{31:03}. Applying the definition of $\xi_y$, the first $d$ components of the image are the coordinates of the base point in the new chart, $(y\circ x^{-1})^i(\alpha^1,\dots,\alpha^d)$, depending only on the $\alpha$'s; the second $d$ components are the components of the vector in the new chart, which are obtained from the old ones by the Jacobian:
\begin{equation}\label{eq:lag-trans}
  \bigl(\xi_y\circ\xi_x^{-1}\bigr)(\alpha,\beta)\;=\;\Bigl(\,(y\circ x^{-1})^i(\alpha)\,,\ \ \frac{\partial(y\circ x^{-1})^i}{\partial\alpha^m}(\alpha)\;\beta^m\Bigr)_{i=1,\dots,d} .
\end{equation}
(The components of a vector in the $y$-system are those of the $x$-system multiplied by $\partial y^i/\partial x^m$: this follows from the definitions in a few lines, which the lecturer ``just writes down'') \ts{13}{1}{31:03}--\ts{13}{1}{33:49}. So the chart transition is expressed completely by the \emph{old} chart transition maps $y\circ x^{-1}$ and their partial derivatives, summed with the components $\beta^m$.

We know that each chart transition $y\circ x^{-1}$ of the smooth atlas $\mathcal A$ is smooth, hence so are all their partial derivatives; and the map \eqref{eq:lag-trans} is built from them just by multiplication with the $\beta^m$ and summation, hence again smooth. This shows: the atlas $\mathcal A_{TM}$ is a smooth atlas, and \emph{$TM$ is a smooth manifold of dimension $2d$} \ts{13}{1}{34:09}--\ts{13}{1}{35:26}.

\begin{theorem}[$TM$ is a smooth manifold of dimension $2d$]\label{thm:lag-smooth}
Let $(M,\mathcal O,\mathcal A)$ be a smooth manifold of dimension $d$. Then the induced atlas $\mathcal A_{TM}$ of \cref{def:lag-atlas} is a smooth atlas of $TM$ (with the topology it induces), so $TM$ is a smooth manifold of dimension $2d$, and $\pi\colon TM\to M$ is smooth.
\end{theorem}
\begin{proof}
The charts cover $TM$; the transitions \eqref{eq:lag-trans} are smooth, as shown above. (The smoothness of $\pi$ is only mentioned in this chapter as ``obvious'': in the charts, $\pi$ is $(\alpha,\beta)\mapsto\alpha$ [my addition, not said in class].)
\end{proof}
The lecturer stresses that everything up to here was a definition, ``and from here on it is hopefully no definition any more but follows from the previous; write it out again at home quietly and then it becomes clear why it must be so'' \ts{13}{1}{35:26}. ``Are there questions? \dots\ This is an absolutely elementary construction that we will need.'' \ts{13}{1}{35:50}.

% ==================================================================
\section{Lifting of curves and the velocity phase space}\label{sec:lag-lift}

What can one imagine physically under this? Before that I need one more concept which in combination with the tangent bundle transports the physical intuition: the \textbf{lifting of curves} \ts{13}{1}{35:50}--\ts{13}{1}{36:31}.

\begin{definition}[New concepts: the lift (Hebung) of a curve]\label{def:lag-lift}
Let $\gamma\colon[a,b]\to M$ be a smooth curve (``in view of our application, in mechanics, with a one-parameter range, a closed interval, so that we can prescribe an initial and an end point''). The \textbf{lift} of $\gamma$ is the curve in the tangent bundle
\begin{equation}\label{eq:lag-lift}
  \tilde\gamma\colon[a,b]\to TM,\qquad \tilde\gamma(\lambda)\;:=\;X_{\gamma,\gamma(\lambda)}\,,
\end{equation}
where $X_{\gamma,\gamma(\lambda)}\in T_{\gamma(\lambda)}M$ is the tangent vector (velocity) of $\gamma$ at the point $\gamma(\lambda)$ (the lecturer writes $\tilde\gamma$ with the symbol of a ``$\gamma_*$'': a curve which does not land in the manifold but in the tangent bundle) \ts{13}{1}{36:31}--\ts{13}{1}{39:21}.
\end{definition}
One knows the tangent vector to a curve at a given point; the new thing is that to each parameter value one assigns \emph{this vector, in the bundle}: a point of the tangent bundle which remembers both where one is and how fast. Note that $\pi\circ\tilde\gamma=\gamma$ \ts{13}{1}{38:38}--\ts{13}{1}{39:21}.

\begin{example}[Lifting a curve in a one-dimensional manifold]\label{ex:lag-liftex}
The lecturer's picture: let $M$ be one-dimensional and imagine you run over it with your curve $\gamma$, starting at $\gamma(a)$ slowly and getting faster and faster until you arrive at $\gamma(b)$. Think of the curve as a trace of luminous points: each second the particle lights up, leaving a dash; short dashes at the beginning, long dashes later. At each point of the manifold attach the tangent space; in the tangent bundle the lifted curve $\tilde\gamma$ passes through the vectors, which are short at the beginning and long at the end.
``On the horizontal axis you read off where you are, and on the vertical axis (which is really only a point in this manifold, one has to identify) how fast you are'': the lift ``unfolds the velocity'' \ts{13}{1}{37:14}--\ts{13}{1}{41:47}.
\end{example}
\begin{aside}[Another example (as said in class)]\label{as:lag-office}
If the manifold is the $S^1$ (``you may paste it together at home from the handout'') a more complicated curve can be drawn: someone leaves his office, gets faster, says ``damn, do I have my pen with me?'', slows down, stops, says ``I must go back to the office, I have the pen in my pocket'', turns around and goes back; this curve is easily shown in the chart of the velocity--position phase space but does not look like much on the circle itself. In the chart one sees the same curve $\xi_x\circ\tilde\gamma$, ``a bit distorted'' \ts{13}{1}{42:08}--\ts{13}{1}{44:31}.
\end{aside}

\paragraph{The chart image of the lift.} Going to the chart: the lift in the chart $\xi_x$ is
\begin{equation}\label{eq:lag-liftchart}
  \xi_x\circ\tilde\gamma(\lambda)\;=\;\Bigl(\gamma^1_{(x)}(\lambda),\dots,\gamma^d_{(x)}(\lambda),\ \dot\gamma^1_{(x)}(\lambda),\dots,\dot\gamma^d_{(x)}(\lambda)\Bigr),\qquad \gamma^i_{(x)}:=x^i\circ\gamma ,
\end{equation}
(the first $d$ components are the chart representation of the curve, the other $d$ its derivatives, the components of the velocity, as one has seen in \cref{eq:lag-comp}) \ts{13}{1}{43:50}--\ts{13}{1}{44:31}. This is what physics knows as the \textbf{velocity phase space} (there is another phase space which we meet later, in the Hamiltonian mechanics): position and velocity. In one dimension it is two-dimensional, for a two-dimensional manifold it is four-dimensional, ``I draw bad pictures, but you do not need pictures, you have the wonderful definition and with that one can calculate everything'' \ts{13}{1}{44:31}--\ts{13}{1}{45:34}.

\begin{coursenote}[Elementary, and how to practise (as said in class)]\label{cn:lag-practise}
``If you go home and look at it again you will find that it is precisely the same stuff that we did in the first weeks, only applied to a new case. We construct an atlas, and at first sight there are many indices again, but it is all child's play.'' If you want to understand it without many indices: take a one-dimensional manifold, put $d=1$, compute all these objects, and draw the pictures \ts{13}{1}{45:34}--\ts{13}{1}{46:15}. Then the lecturer invites questions (``they need not be technical, you may also ask: what is the idea again?'') and asks who has the feeling of having understood almost nothing \ts{13}{1}{46:15}--\ts{13}{1}{47:00}.
\end{coursenote}

\begin{intuition}[Why the tangent bundle? The lift as an alternative to the spacetime picture]\label{int:lag-why}
Think of the manifold as physical space and of curves in it as particle trajectories. If the particle blinks, one sees long and short dashes and can estimate the speed; but if it shines all the time one only has a \emph{point set}: from the set alone one cannot know where the particle was fast or slow, and that is essential information. One way of making it visible beyond blinking is what we used all the time: a \emph{spacetime picture}: one resolves the spatial motion in spacetime (of dimension $d+1$), and a faster curve is steeper.
Here we do something more dramatic: we go to the \emph{double dimension} and afford the luxury of recording, over each point where the particle is, the full vector (for which one needs $d$ more components); so we get a $2d$-dimensional velocity phase space. As always we have a definition in the abstract (``in the real world'') and one can map it to a chart; if you look at a book (e.g.\ Flie{\ss}bach) you get everything served in a chart, ``and there it is written so simply: why is it not explained like that in the lecture?'' \ts{13}{1}{47:00}--\ts{13}{1}{49:04}.
\end{intuition}
The point is that one can calculate with it, but then one cannot ask conceptual questions, for instance: what is the \emph{canonical momentum}, what kind of object is it, a vector or a covector? This cannot be answered in a chart: one must look at the real world, and without the technology one cannot learn anything about it. ``It is incredibly important that we shall find that the momentum, as a mathematical object, is \emph{not a vector}: it will turn out to be a \emph{covector}.'' That one can understand only if one understands this mathematics \ts{13}{1}{49:04}--\ts{13}{1}{50:06}.

\begin{coursenote}[Break, and a greeting (as said in class)]\label{cn:lag-break}
Break until 11 o'clock, then it goes on. An important announcement: ``you all greet so nicely with `Mechanik', but these could also be some who were there only for the first time; so from now on greet with `Phasenraum' (phase space)''; the lecturer wonders who will try it \ts{13}{1}{49:45}--\ts{13}{1}{50:27}.
\end{coursenote}

% ==================================================================
\section{The Lagrangian}\label{sec:lag-L}

Now one of the central objects of classical mechanics, found in every textbook: the Lagrangian \ts{13}{1}{50:27}--\ts{13}{1}{50:47}.

\begin{definition}[New concepts: Lagrangian (Lagrange-Funktion)]\label{def:lag-L}
A \textbf{Lagrangian} is a smooth function
\begin{equation}\label{eq:lag-L}
  L\colon TM\to\R .
\end{equation}
It lives on the tangent bundle (the velocity phase space): to every point of the velocity phase space, that is, to every point of your handout, it assigns a real number. There are \emph{no further conditions} \ts{13}{1}{50:47}--\ts{13}{1}{51:49}.
\end{definition}
``So the central object of mechanics is something very simple: a smooth function on the velocity phase space, or the tangent bundle of the space on which our system lives. That is it; so simple. You see: it pays to define these things.'' \ts{13}{1}{51:49}.

\begin{coursenote}[The Lagrangian and charts (as said in class)]\label{cn:lag-chartfree}
What does it have to do with charts? Nothing: one need not choose any chart to write down the Lagrangian; it is an object living in the real world. Whichever chart one prefers (``I want this chart or that chart or this atlas''), the Lagrangian is ``completely cold'' about it. In a concrete case one \emph{specifies} the Lagrangian by saying ``I choose my favourite chart and specify it there''; but it is important to realise that the object has first nothing to do with a chart \ts{13}{1}{52:09}--\ts{13}{1}{52:51}.
\end{coursenote}

Now I may \emph{wilfully} choose a chart $(U,x)$ of the base manifold's atlas; it generates $TU\subseteq TM$ and the induced chart map $\xi_x\colon TU\to\R^{2d}$ of \cref{def:lag-xi}. On $TU$ the Lagrangian is still defined (restriction $L|_{TU}$) and we can write it in terms of numbers:

\begin{definition}[New concepts: chart representation of the Lagrangian]\label{def:lag-Lchart}
\begin{equation}\label{eq:lag-Lchart}
  L_{(x)}\;:=\;L\circ\xi_x^{-1}\colon\ \xi_x(TU)\subseteq\R^{2d}\to\R ,
\end{equation}
the \textbf{chart representation of the Lagrangian}, a function of $2d$ real numbers (cf.\ \cref{def:dm-fnrep}): one takes a $2d$-tuple, goes up with $\xi_x^{-1}$ to the tangent bundle and then applies $L$ \ts{13}{1}{52:51}--\ts{13}{1}{54:36}.
\end{definition}
``What you find in mechanics books as Lagrangians is always this chart representation; this is not said explicitly, but it is important to understand it: if one does not, one wonders what it means to change coordinates and charts in these books, what the people are talking about, since the Lagrangian is such and such. Not at all: what they write is the \emph{representation in a particular chart}, and above it, in the real world, there stands the true object that interests us'' \ts{13}{1}{54:36}--\ts{13}{1}{55:38}.

\begin{example}[The Lagrangian of the harmonic oscillator, in a chart]\label{ex:lag-ho}
Take a one-dimensional manifold, no matter which one, and some chart (``which nobody in the textbook tells you, there it says suddenly ``with respect to a chart that may exist'''): the chart representation $L_{(x)}$ depends on $2d=2$ real numbers, call them $(\alpha,\beta)$ (position and velocity component; ``the book would write $q$ and $\dot q$ or $x$ and $\dot x$''). In the textbook one finds then, for the harmonic oscillator,
\begin{equation}\label{eq:lag-ho}
  L_{(x)}(\alpha,\beta)\;=\;\tfrac12\,m\,\beta^2-\tfrac12\,k\,\alpha^2 ,
\end{equation}
the kinetic energy minus the potential energy of the harmonic oscillator, with $m$ the mass and $k$ the spring constant (the lecturer writes first ``$\tfrac12\beta^2-\dots$'' and wonders whether to put the $\tfrac12$ and the minus sign; ``I am on the safe side with the $\tfrac12$''; the captions are garbled [?]) \ts{13}{1}{55:38}--\ts{13}{1}{58:25}. In books one reads instead $\tfrac12 m\dot x^2-\tfrac12kx^2$ and says that $\dot x$ is the derivative of $x$; but for this to have a meaning one must realise that these are \emph{just points of a $2d$-dimensional space}, not curves. ``The first thing to know about Lagrangians is that they depend on $2d$ variables, and these stand there as such.''
Why this is the Lagrangian of the harmonic oscillator (``which has until now nothing to do with the oscillator'') will be explained later; ``you have computed it all on your problem sheets in a philosophical way'', and I will explain what it is good for in the next lecture; first one must understand that the object is abstract, and what it is not: it is, in a chart, just such a function \ts{13}{1}{56:00}--\ts{13}{1}{58:46}.
\end{example}
One usually takes a chart in order to put a concrete function on the table in an exam; so ``what is written there'' \ts{13}{1}{58:46}.

\paragraph{Partial derivatives and notation.} Now this is a function of $2d$ variables, so one can differentiate it in the chart. For the example: differentiate with respect to the first entry and evaluate at $(\alpha,\beta)$: one gets $-k\alpha$; differentiate with respect to the second entry: $m\beta$ \ts{13}{1}{58:46}--\ts{13}{1}{59:48}. More generally, for a chart representation $L_{(x)}\colon\xi_x(TU)\subseteq\R^{2d}\to\R$ we can form all the partial derivatives with respect to the $a$-th entry, $\partial_aL_{(x)}$; in applications the first $d$ entries (positions) are treated differently from the last $d$ (velocities), so one writes:
\begin{equation}\label{eq:lag-partials}
  \partial_aL_{(x)}\ (a=1,\dots,d)\ \text{ and }\ \partial_{d+a}L_{(x)}\ (a=1,\dots,d),
\end{equation}
that is, with $a=1,\dots,d$ one differentiates with respect to the first $d$ entries (the positions), and with $d+a$ with respect to the last $d$ (the velocity components); this is what we need in the calculation \ts{13}{1}{59:48}--\ts{13}{1}{1:01:58}. One could put all of them from $1$ to $2d$ over one comb; but in the applications it is usually the case that they are treated differently, and then it is clever to split them.

\begin{coursenote}[The traditional notation for these derivatives, and why it is not used here (as said in class)]\label{cn:lag-notation}
A mysterious notation that is ``actually customary but not practised by us'': one invents a new symbol, $\partial L/\partial q^a$ for the first and $\partial L/\partial\dot q^a$ for the second family (``this is not on the manifold, it is here in the chart''). It is nicer, up to the point where one inserts something, e.g.\ $\dot q$, as the component of a tangent vector; then one suddenly differentiates with respect to $\dot q$ and asks what this means. This is introduced in every mechanics lecture, and then one is asked: ``why can I differentiate with respect to $\dot q$ if I differentiate with respect to $q$, what does it mean?'' And it is always said: ``one treats them as independent''.
But these are mathematical objects: ``that you can treat them so and so, once they scream and once they do not''; the only mathematically clear and consistent notation is the one written above; the other leads sooner or later to confusion; hence we do not use it. But you should know it if you want to compare with a textbook \ts{13}{1}{1:01:58}--\ts{13}{1}{1:04:49}.
\end{coursenote}

\paragraph{Level lines of the harmonic oscillator.} Since it is a function on the velocity phase space, I can draw in the chart $(\alpha,\beta)$ the \emph{level lines} of this function: a function on a two-dimensional plane; one wants to know what it looks like \ts{13}{1}{1:04:49}--\ts{13}{1}{1:05:54}. In the chart one may think of $\beta$ (vertical) as the velocity component and $\alpha$ (horizontal) as the position component; ``but it need not be, first it is just a function''. For the typical theorist's choice $k=m=1$, $L_{(x)}(\alpha,\beta)=\tfrac12(\beta^2-\alpha^2)$ vanishes exactly on the cross $|\alpha|=|\beta|$; the other level lines are hyperbolas, on both sides of the cross, with the values $\pm$ some positive number, $\pm 2,\dots$, symmetrically (the lecturer reads off $0$, $\pm1$, $\pm2$ and so on from the sketch; the captions are garbled in this passage [?]) \ts{13}{1}{1:05:54}--\ts{13}{1}{1:08:28}.

\begin{intuition}[What the Lagrangian replaces]\label{int:lag-replaces}
``You cannot yet understand it, I have not shown it yet'': in the Newtonian formulation of the harmonic oscillator one specifies the force law, the right-hand side of $m\ddot x=\dots$; the left-hand side is always the same, and different physical systems differ in the Newtonian formulation by different specifications of concrete forces and their dependence on position, and so on. In the Lagrangian mechanics it turns out that \emph{everything one needs to know about the harmonic oscillator, in particular the force law, is already encoded in such a function}, and precisely this function is that of the harmonic oscillator.
If I give you another system (recall the funny thing with the pendulums and springs at the beginning), I give you another such function; take the same chart and give a rather complicated physical system: springs that are weaker and stronger in some regions, and who knows what; if you have such a system in the laboratory and measure what it does, you can in principle construct the Lagrangian from practice, and then comes the question: all very well, but the equation of motion cannot be solved by anybody; \dots\ \ts{13}{1}{1:08:28}--\ts{13}{1}{1:11:00}.
\end{intuition}

\begin{aside}[Outlook: the Hamiltonian formulation (as said in class)]\label{as:lag-hamilton}
A short outlook: this was the velocity phase space, where we have something like position and velocity components, but later we shall get to know \emph{another phase space}, that of the \textbf{Hamiltonian mechanics}. There we have again a phase space, but it will not be the tangent bundle: it will be the \emph{cotangent bundle}: to each point one attaches the cotangent space, and in the chart in simple cases one axis is the position and the other one is a \emph{momentum}. Under certain conditions one can turn a Lagrangian into a \textbf{Hamiltonian} (a capital $H$, not an $L$), which lives not on the tangent bundle but on the cotangent bundle and is smooth there; we can convert them into each other, and this is a different representation.
The lecturer writes the Hamilton function that corresponds to the Lagrangian of the oscillator (not understood yet) and notes: for the ``terrible buckling system'' the level lines would be much more complicated; but in Hamiltonian mechanics he will present a \emph{geometric procedure} with which one can discuss systems that nobody can calculate: via a so-called \emph{symplectic structure} one constructs a \emph{Hamiltonian vector field}, and the system follows the \emph{integral curves} of this vector field; one can then draw them, or write them with the computer. ``A great method to predict how it looks like for things which you cannot calculate by hand any more; it is child's play, we only have to develop some mathematics, so that you see it is correct.'' But first we deal with Lagrangians for a long time, ``since with these one can already do damn much'' \ts{13}{1}{1:11:00}--\ts{13}{1}{1:13:48}.
\end{aside}

% ==================================================================
\section{The action functional}\label{sec:lag-action}

You now say: all very well, you give me a Lagrangian and claim it has something to do with a physical system, what is all this? We begin to tie it together, and the decisive object is the \textbf{action functional (Wirkungsfunktional)} \ts{13}{1}{1:13:48}--\ts{13}{1}{1:14:39}.

\begin{definition}[New concepts: the action functional (Wirkungsfunktional) belonging to a Lagrangian]\label{def:lag-action}
Let $L\colon TM\to\R$ be a Lagrangian (a smooth function on the tangent bundle). Let $p,q\in M$ and
\begin{equation}\label{eq:lag-curves}
  \mathcal C^\infty_{p,q}\;:=\;\bigl\{\gamma\colon[a,b]\to M\ \text{smooth}\ \big|\ \gamma(a)=p,\ \gamma(b)=q\bigr\}
\end{equation}
the set of all smooth curves from the point $p$ to the point $q$ (the lecturer writes $C^\infty([a,b],M)$ with $p$ and $q$ as decoration; $p,q$ may be chosen as one likes, but for a concrete problem one chooses a pair). The \textbf{action functional} belonging to $L$ is the map
\begin{equation}\label{eq:lag-S}
  S\colon\mathcal C^\infty_{p,q}\to\R,\qquad S[\gamma]\;:=\;\int_a^b\mathrm d\lambda\ L\bigl(\tilde\gamma(\lambda)\bigr),
\end{equation}
where $\tilde\gamma$ is the lift of $\gamma$ (\cref{def:lag-lift}). The real number $S[\gamma]$ is called the \textbf{action} (Wirkung) of the curve $\gamma$ \ts{13}{1}{1:14:39}--\ts{13}{1}{1:24:33}.
\end{definition}

\begin{coursenote}[Functional versus function; the square brackets; a record book example (as said in class)]\label{cn:lag-functional}
What does an action functional eat? An entire \emph{curve}, and it gives out a real number: the domain is the set of all smooth curves from $p$ to $q$ in the manifold; this is why it is called a functional and not a function: ``it is a kind of function, a map; but a functional''. A map that is not of the form \eqref{eq:lag-S} is just as good a functional: for instance ``you give me a curve and the functional is the highest speed you reached during the whole curve'': the functional does not care how fast you were except for the maximum; ``for the Guinness book of records, the man ran at some time 80 km per hour, and he qualifies''. Evidently we shall not consider such general functionals but only those given by \eqref{eq:lag-S} \ts{13}{1}{1:15:45}--\ts{13}{1}{1:18:55}.
The square brackets, $S[\gamma]$ instead of $S(\gamma)$, are a tradition: to indicate, for whoever slept through the definition, that what stands inside is no number on which the map depends but a whole map; it means nothing otherwise \ts{13}{1}{1:18:55}--\ts{13}{1}{1:19:44}.
\end{coursenote}
Look at the integrand: $L$ is defined on $TM$, and $\tilde\gamma(\lambda)\in TM$; so $\lambda\mapsto L(\tilde\gamma(\lambda))$ is a real-valued function of a real variable on $[a,b]$, which we can \emph{integrate} from $a$ to $b$: ``a child's play, when one understands the Lagrangian like this and takes the lifted curve''. Everything on the right-hand side is either the Lagrangian that I prescribe or the lift of $\gamma$ \ts{13}{1}{1:19:44}--\ts{13}{1}{1:22:30}.

\begin{intuition}[Picture: the action as an integral of the Lagrangian along the lifted curve]\label{int:lag-actionpic}
In the chart picture: the Lagrangian is given as it is, with its level lines, and describes the system (e.g.\ the level lines of an oscillator; and a concrete one for a concrete system). Now there is a curve $\gamma$ which runs from $p$ to $q$, with some speed profile (slow at the beginning, then faster, then slower), so that the lift $\tilde\gamma$ is a certain line in the $(\alpha,\beta)$-plane (in the lecturer's drawing the red $\xi_x\circ\tilde\gamma$; a different curve gives another path). We do not care how the Lagrangian looks in the plane: we only care what the value of the Lagrangian is when I wander \emph{along this curve}; this is a function of $\lambda$ and I integrate over it, so I get a number: for each curve a number, built from the Lagrangian of the lifted curve. This is called the \textbf{action} $S[\gamma]$ of the curve $\gamma$ \ts{13}{1}{1:22:30}--\ts{13}{1}{1:24:33}.
\end{intuition}

\begin{aside}[Feynman and the path integral (as said in class, not checked)]\label{as:lag-feynman}
``I tell you right away what it is about in mechanics, but the action is a quantity that has a meaning, even a central one, up to quantum mechanics'': Feynman, ``one of the greatest of the last century in theoretical physics'', showed that quantum mechanics, which one normally writes with the Schr\"odinger equation, can be written as a quantisation of a classical system by using directly this action. You may have heard that in quantum mechanics a particle thrown at someone's head has not taken one path but \emph{all conceivable paths}, a superposition; but the point is: for each of these conceivable paths one computes the action, and then one ``mashes'' these actions in a so-called \emph{path integral}.
The classical mechanical action stands there, unchanged; if one computes the path integral one gets the transition amplitude, which is complex, and from it the probability that my chalk hit the head of the other person. So there is a formulation of quantum mechanics that uses exactly this action; that is why the action functional is so important. But in the end it is not very difficult: one gives functions, which one has to model, as one had to model the force law \ts{13}{1}{1:24:33}--\ts{13}{1}{1:26:44}.
\end{aside}

\subsection{The action in a chart}\label{ssec:lag-actionchart}

The whole can be written in a chart; one should do it, for two reasons: in the introductory textbooks one sees only the chart representation, and one must know where it comes from; and if I must put something concrete on the table I must always give it in some chart \ts{13}{1}{1:26:44}--\ts{13}{1}{1:27:27}. We choose a chart $(U,x)$; with it we can understand only a part of the manifold, on $TU$ the restriction of the Lagrangian; the lifted curve may not stay in $TU$: the preimage of $TU$ under the lift lies in the interval $[a,b]$ but may be smaller (``it could be fewer'') \ts{13}{1}{1:27:27}--\ts{13}{1}{1:28:09}. We have the induced chart $\xi_x\colon TU\to\R^{2d}$ and the chart representation $L_{(x)}=L\circ\xi_x^{-1}$.

Now insert a chart wilfully (``in earlier days one said that a gentleman chooses a chart only when he really has to, and today a lady also''): in
$S[\gamma]=\int_a^b\mathrm d\lambda\,L(\tilde\gamma(\lambda))$ insert the identity $\xi_x^{-1}\circ\xi_x=\mathrm{id}$, which does nothing, but we can read it diagrammatically as going down to the chart and up again (``this is the same: one has gone down and up; one may insert it thoughtlessly in the diagram, as long as one knows that the successive application is associative, so that one may write it this way''):
\begin{equation}\label{eq:lag-Schart}
  S[\gamma]\;=\;\int_a^b\mathrm d\lambda\ \bigl(L\circ\xi_x^{-1}\bigr)\bigl(\xi_x\circ\tilde\gamma(\lambda)\bigr)\;=\;\int_a^b\mathrm d\lambda\ L_{(x)}\Bigl(\gamma^1_{(x)}(\lambda),\dots,\gamma^d_{(x)}(\lambda),\ \dot\gamma^1_{(x)}(\lambda),\dots,\dot\gamma^d_{(x)}(\lambda)\Bigr),
\end{equation}
where by \eqref{eq:lag-liftchart} the lifted curve in the chart is the $2d$-tuple of the components of the curve in the $x$-chart, $\gamma^i_{(x)}=x^i\circ\gamma$, and their derivatives with respect to $\lambda$ (the lecturer first writes the components of the lift as ``the $x^1$ of the lifted curve'' and then converts them: ``this is the derivative with respect to $\lambda$ of the chart representation of the curve evaluated at $\lambda$'') \ts{13}{1}{1:28:09}--\ts{13}{1}{1:33:24}.

If you now suppress all the decorations (forget that a chart was chosen and that all came from a curve in the real world) you can write, with $q^i(\lambda):=\gamma^i_{(x)}(\lambda)$:
\begin{equation}\label{eq:lag-Sbook}
  S[q]\;=\;\int_a^b\mathrm d\lambda\ L\bigl(q^1(\lambda),\dots,q^d(\lambda),\dot q^1(\lambda),\dots,\dot q^d(\lambda)\bigr),
\end{equation}
and this is what you find in Flie{\ss}bach and elsewhere (``textbook, page 100 or so''): the action functional of a curve is the integral over the Lagrangian of the coordinates where the curve passes and of the derivatives of these component functions. It is all only in the chart picture; and the great disadvantage is that one often does not notice that a chart has been chosen: in another chart all these objects have another form, looking just the same but being other objects; so if one suppresses the information that it holds only in a particular chart, much of what follows becomes conceptually mysterious. ``We are not mysterious; ours is a bit more complicated, but in theory more complicated is always better than mysterious: we can deal with complexity, but not with the mysterious'' \ts{13}{1}{1:33:24}--\ts{13}{1}{1:35:30}.

\begin{coursenote}[Warning on the formula \eqref{eq:lag-Sbook}]\label{cn:lag-warn}
(My remark, not said in class.) The abbreviation $q^i$ for $\gamma^i_{(x)}$ and $\dot q^i$ for its derivative hides that in \eqref{eq:lag-Sbook} $L$ is the chart representation $L_{(x)}$, and that $(q,\dot q)$ are the chart coordinates of the lifted curve, not independent variables unrelated to $\gamma$.
\end{coursenote}

% ==================================================================
\section{Goal and strategy of the calculus of variations}\label{sec:lag-strategy}

Now you have a Lagrangian and can build an action functional; all very well, but what does it have to do with mechanics, and above all with Newton's second axiom? This is still obscure. I did it deliberately like this, so that we first understand the objects mathematically and only then give them physical meaning: otherwise one gives physical meaning too quickly, and at some point one has interpreted more into it than is really there, and then it becomes problematic. ``I ask you, as so often, to accept this delayed gratification'' \ts{13}{1}{1:35:30}--\ts{13}{1}{1:36:33}.

So for the last 15 minutes: \textbf{goal and strategy of the calculus of variations (Variationsrechnung)}. Next time we shall find that we can write Newton's second axiom as a \emph{variational principle}: one demands that, for this action functional, a certain quantity, its \emph{variation}, be zero. All of the information is in the action functional; from it we shall simply pull out the second Newtonian law by the calculus of variations. ``It just comes out'', and then the connection is established \ts{13}{1}{1:36:33}--\ts{13}{1}{1:37:14}.

\begin{intuition}[Systems compose by adding actions]\label{int:lag-sum}
I tell you already (before I say again what the calculus of variations is): write down the action for a pendulum; take the action for another pendulum; the action of the combined system of the two pendulums is, and this is the super new message, \emph{the sum of the original actions}. If you now connect the two pendulum masses by a spring, the new action of the whole system, which really describes it, is
\[
  S_{\text{total}}\;=\;S_{\text{pendulum 1}}+S_{\text{pendulum 2}}+S_{\text{interaction}},
\]
with a so-called interaction term, which is also a functional: both the objects themselves and the interaction through a spring, or through whatever fields, all are represented as a \emph{sum of actions}. Composing systems means just adding actions \ts{13}{1}{1:37:14}--\ts{13}{1}{1:38:16}.
Imagine the complicated system of the beginning, with all the rods and the beads: try to decompose it into forces and write down the force law: ``almost not feasible''. Instead you analyse the system nicely piece by piece, determine the action of each piece, add all up, and then you have the action for the whole system; then you carry out the \emph{calculus of variations}, and out comes a gigantic complicated equation of motion, differential equations; in general you cannot solve them by hand if the system is too complicated, but they are the \emph{right equations}: you know at least what you want to solve \ts{13}{1}{1:38:16}--\ts{13}{1}{1:38:57}.
Moreover, if you restrict parts of the system to subspaces (the bead to the wire), this is called a \emph{constraint}, and ``we also get that for free with the action functional; it is incredibly simple'' \ts{13}{1}{1:43:05}--\ts{13}{1}{1:43:46}.
\end{intuition}

\begin{secondary}[Beyond the core: how physics works, and the ``dirty arrow'' of modelling (as said in class)]\label{sec:lag-modelling}
This is very important: with it we have turned a physical problem of great complexity into a \emph{purely mathematical problem}, the solution of a differential equation: ``this is super theory''. But one must see how physics works: one has a problem in the real world (``in the real world''), which one converts into mathematics, depending on how naive one is, by writing down something like an equation; and \emph{this step requires physical feeling}, it requires the physicist; a pure mathematician who has just left his office knows little of the world. This step is a \emph{dirty arrow}, for one does all kinds of terrible things: one models the length of the pendulum rod by a real number, and this recurs again and again; it is absorbed in the Lagrangian.
Consider that if one took the modelling too seriously, a pendulum rod would have a real length, and I could realise every real length in nature, whereas in reality one can only produce multiples of atomic diameters or the like as lengths. A modelling always makes idealisations that one must not take too seriously; one may calculate with them, but must always take care that everything is still all right. This arrow is called \textbf{modelling}: how do I turn a physical problem into a mathematical one \ts{13}{1}{1:38:57}--\ts{13}{1}{1:41:43}.
Then come our friends the mathematicians, and sometimes we ourselves: we use mathematics and derive, e.g., $x(t)=A\cos(\omega t+\varphi_0)$ as the solution of a differential equation. Then comes again a dirty step: the \emph{interpretation} of what one has actually computed: ``you said you could model the length as a real number, and now you interpret back: I can realise any length in nature; it was not meant like this''; everything is dirty there, and one has to take care. Then you let the time run, and suddenly the pendulum is there (the experimental physicist simply says: I look at what happens; this is the left side) and the two must somehow agree.
But one must always see: the mathematics in the middle is exact and super and important, while in the modelling one can produce a lot of nonsense or dirt, in the interpretation too, and here one needs a physicist \ts{13}{1}{1:41:43}--\ts{13}{1}{1:42:44}.
\end{secondary}
So we say: if we can model the small systems, then we can handle the big complicated systems with springs and rods simply by adding the actions; and this is why we want to learn how to obtain Newton's law from the action: with these action functionals one can put systems together incredibly easily, and also (as I have shown, if one restricts parts of the system to subspaces, such as the bead on the wire, a \emph{constraint}) get constraints ``for free'' \ts{13}{1}{1:42:44}--\ts{13}{1}{1:43:46}.

\begin{coursenote}[End of the lecture]\label{cn:lag-end}
The lecturer decides to leave the full goal and strategy of the calculus of variations for next time: ``at the end next time stand the Euler--Lagrange equations, and then you see the hundred percent contact with Newtonian theory; we have the reformulation. Until next time'' \ts{13}{1}{1:43:46}.
\end{coursenote}

% ==================================================================
\section*{Key concepts}
\begin{keyconcepts}
\textbf{Retrieval questions.}
\begin{enumerate}
  \item Define the tangent bundle $TM$ and the bundle projection $\pi$. What is $TS^1$ as a set?
  \item Construct the induced chart $\xi_x$ on $TU$. What are its first $d$ and its last $d$ components, and why is the dimension $2d$?
  \item Write the transition map $\xi_y\circ\xi_x^{-1}$ and use it to show that the induced atlas is smooth.
  \item Define the lift of a curve; what is its image in a chart? What does the lift record that the point set of the curve does not?
  \item Define a Lagrangian. Why does it not need a chart? What is its chart representation, and what do textbooks mean by $L(q,\dot q)$?
  \item Why is the notation $\partial L/\partial\dot q$ problematic, and how do we write these derivatives?
  \item Define the action functional of a Lagrangian: its domain, its value on a curve; why is it a functional and not a function?
  \item Which facts about actions will make the Lagrangian formulation convenient (sum of actions, constraints), and what is modelling versus interpretation?
\end{enumerate}
\textbf{Concepts.}
\begin{itemize}
  \item Tangent bundle and bundle projection (\cref{def:lag-tb}, \eqref{eq:lag-tb}, \eqref{eq:lag-pi}); the induced atlas and $TM$ as a $2d$-dimensional manifold (\cref{def:lag-xi,def:lag-atlas,thm:lag-smooth}, \eqref{eq:lag-xi}, \eqref{eq:lag-trans}).
  \item Lift of a curve, velocity phase space (\cref{def:lag-lift}, \eqref{eq:lag-lift}, \eqref{eq:lag-liftchart}).
  \item Lagrangian and its chart representation; harmonic oscillator; partial derivatives (\cref{def:lag-L,def:lag-Lchart,ex:lag-ho}, \eqref{eq:lag-L}, \eqref{eq:lag-Lchart}, \eqref{eq:lag-ho}, \eqref{eq:lag-partials}).
  \item Action functional and its chart form (\cref{def:lag-action}, \eqref{eq:lag-S}, \eqref{eq:lag-Schart}, \eqref{eq:lag-Sbook}).
  \item Outlook: Euler--Lagrange equations, Hamiltonian formulation on the cotangent bundle (\cref{as:lag-hamilton,int:lag-sum}).
\end{itemize}
\end{keyconcepts}

\flushside
\end{paracol}
